\section{Open-Weight 研究窗口：解释、攻击与防御}

\subsection{可见内部带来的研究机会}
进入 open-weight 层后，研究者可做的事情显著增加：表征可解释、激活分析、特征编辑、蒸馏与对齐实验都成为可能。

\begin{importantbox}{Open Weights 带来的是“机制研究权”，不是“万能改造权”}
你可以深入模型内部，但仍被训练历史、架构形态和数据来源所约束。  
因此，open-weight 是强能力层，但不是终局层。
\end{importantbox}

\subsection{安全侧：能力提升与风险暴露同步发生}
讲座中段多次提到安全相关实验：开放权重让防御研究更深入，也让攻击路径更可操作。  
这不是“开或不开”的单选题，而是“如何配套治理与审计机制”的系统题。

\begin{table}[H]
\centering
\begin{tabular}{p{2.8cm}p{4.0cm}p{4.8cm}}
\toprule
\textbf{研究方向} & \textbf{开放权重的收益} & \textbf{对应风险} \\
\midrule
机制解释 & 可定位失败模式与脆弱特征 & 攻击者也可定位薄弱环节 \\
对齐实验 & 可复现训练与安全干预路径 & 低成本复制有害变体 \\
防御评测 & 可进行高强度红队与压力测试 & 防御细节被对抗性学习利用 \\
\bottomrule
\end{tabular}
\caption{Open-weight 研究中的“能力-风险同增”结构}
\end{table}

\begin{warningbox}{技术社区常见误判：把风险治理外包给“模型提供方”}
一旦生态进入开源/开权重形态，风险治理不能只靠单一组织。  
需要基准、审计工具、部署规范与责任机制共同演进。
\end{warningbox}

\subsection{本章小结}
Open-weight 打开了机制研究和安全研究的重要窗口，但它同时放大了生态治理难度。技术进步必须与治理能力同步扩容。


